\chapter{Adversarial Robustness and Active Probing}
\label{ch:adversarial-robustness}

This chapter keeps the adversarial story grounded in the battery artifacts
that are actually locked in the repo.  The current evidence package is a
spoofing-detection audit rather than a full controller-versus-controller
competition under malicious sensing.

\section{Battery Threat Model}

The battery-specific adversary considered here tries to make the controller
trust a telemetry stream that looks statistically healthy while drifting away
from the true plant state.  In practice that means:

\begin{enumerate}
  \item replaying or smoothing readings so the battery appears stable,
  \item spoofing load-sensitive inputs so the dispatch policy overcommits, and
  \item hiding the attack inside otherwise plausible packet behavior.
\end{enumerate}

The point of active probing is to break that illusion by asking the plant for
a small physical response that a forged telemetry stream cannot mimic forever.

\section{Locked Detection Package}

The current battery outputs are:

\begin{itemize}
  \item \texttt{reports/probing/spoofing\_attack\_results.csv}, which stores
    the attack-detection confusion counts, and
  \item \texttt{reports/probing/probing\_detection\_quality.png}, which turns
    the main detection metrics into a chapter-facing figure.
\end{itemize}

These are generated from the underlying publication artifact
\path{reports/publication/active_probing_spoofing_detection.csv}.

\section{Detection Results}

The locked spoofing-detection audit reports 300 evaluated events with a
threshold of 0.6, producing:

\begin{itemize}
  \item 48 true positives,
  \item 6 false positives,
  \item 4 false negatives,
  \item 242 true negatives,
  \item 0.889 precision, and
  \item 0.923 recall.
\end{itemize}

\begin{table}[ht]
\centering
\caption{Locked active-probing detection results for the battery threat model.}
\label{tab:spoof-results}
\begin{tabular}{lr}
\toprule
\textbf{Metric} & \textbf{Value} \\
\midrule
Events evaluated & 300 \\
Threshold & 0.60 \\
True positives & 48 \\
False positives & 6 \\
False negatives & 4 \\
True negatives & 242 \\
Precision & 0.889 \\
Recall & 0.923 \\
\bottomrule
\end{tabular}
\end{table}

\section{Battery-Operational Meaning}

For the battery book, these numbers mean the probing layer is doing useful
work: it catches most spoofing events while keeping false alarms bounded.
That matters because every false alarm tightens the battery unnecessarily,
while every missed detection risks preserving a hidden observation-action
safety gap.

\section{Bounded Claim}

This chapter now makes the narrower and more honest claim supported by the
locked artifacts:

\begin{itemize}
  \item The battery implementation includes an active-probing detection path.
  \item That detection path achieves high recall and strong precision on the
    current spoofing audit.
  \item An adversarial OQE mode (\texttt{adversarial\_mode} in
    \texttt{dc3s.yaml}) now supplements detection with a real-time
    Byzantine-robust reliability estimator.  When enabled, the runtime takes
    $w_t = \min(w_t^{\mathrm{heuristic}},\, w_t^{\mathrm{robust}})$,
    ensuring that an attacker who can influence individual sub-scores
    cannot raise the aggregate score above the robust floor computed over
    a sliding window of recent signal readings via
    \texttt{compute\_reliability\_robust()}.
  \item The chapter does \emph{not} yet prove adversarial completeness or a
    full safety frontier against all attack strategies; game-theoretic
    optimality remains open, and the robust estimator provides a bounded
    defence rather than a provably optimal one.
\end{itemize}

\section{Summary}

The adversarial extension is no longer just conceptual.  The repo now
contains a locked battery spoofing-detection package with explicit confusion
counts, precision, recall, and a chapter-facing figure.  That gives the
battery thesis a real active-probing evidence surface while keeping the claim
appropriately bounded to what has actually been implemented and measured.
